% !TeX root = ../../Book3.tex
% 纲要标题：### 25.2 规范固定与椭圆化
% 纲要位置：几何学计划纲要.md，第 3176 行。
% 正文状态：初稿；尚未正式审核。

% 纲要内容开始
% 
% * 散度、Bianchi 算子和规范条件
% * 规范化 Einstein 算子的主符号
% * 切片所需的椭圆方程
% * 无穷小等距与剩余对称
% 
% 纲要内容结束

\section{规范固定与椭圆化}
\label{g3:sec:2502}

同一度量经过坐标重参数化，会给出无穷多个张量表达。若直接对 Ricci 方程求逆，这些表达沿微分同胚轨道产生的方向便落在线性化的核中。规范固定要在每条附近的轨道上选一个代表，再把这个代表条件加入方程。这里使用散度条件，并保留 Einstein 张量的散度恒等式；这样既能得到椭圆算子，也能证明新增的规范项在解上消失。

本节固定闭连通光滑流形 \(M^n\)、\(n\ge3\)，以及光滑 Einstein 度量 \(g_0\)，满足 \(\operatorname{Ric}_{g_0}=\lambda g_0\)。下标 \(0\) 表示用 \(g_0\) 定义的微分、内积与体积密度。所有 Sobolev 范数均以这个背景计算。整数 \(s>n/2+3\) 将保证所用张量至少有三阶连续导数。

\subsection{散度、Bianchi 算子和规范条件}
\label{g3:subsec:2502:01}
回顾上一节的约定
\[
	(\delta_g h)_j=-\nabla^i h_{ij},\qquad
	(\delta_g^*\omega)_{ij}=\tfrac12(\nabla_i\omega_j+\nabla_j\omega_i),\qquad
	\beta_g h=\delta_g h+\tfrac12\mathrm d\operatorname{tr}_g h.
\]
闭性使 \(\delta_g^*\) 确为 \(\delta_g\) 的 \(L^2(g)\) 伴随。令 \(X^\flat=g(X,\cdot)\)，则 \(\mathcal L_Xg=2\delta_g^*X^\flat\)。因此散度为零的张量与所有无穷小重参数化正交。我们采用的代表条件为
\begin{equation}\label{g3:eq:metric-divergence-gauge}
	\delta_0(g-g_0)=0.
\end{equation}
这里固定的是背景散度；不能把它换成 \(\delta_g g=0\)，后者对每个度量都自动成立。

\begin{proposition}\label{g3:prop:bianchi-gauge-symbol}
	设 \((M^n,g)\) 为光滑黎曼流形，\(n\ge2\)，\(\omega\) 为光滑一形式。则
	\begin{equation}\label{g3:eq:bianchi-killing-operator}
		2\beta_g\delta_g^*\omega
		=\nabla^*\nabla\omega-\operatorname{Ric}_g(\omega,\cdot).
	\end{equation}
	算子 \(2\delta_g\delta_g^*\) 的主符号为
	\[
		\omega\longmapsto |\xi|_g^2\omega+\xi\langle\xi,\omega\rangle_g,
		\qquad \xi\ne0,
	\]
	因而强椭圆。若 \(M\) 闭，则其核恰为满足 \(\delta_g^*\omega=0\) 的一形式。
\end{proposition}
\begin{proof}
	在一点取正交法标架。由协变导数交换公式，
	\[
		(2\delta_g\delta_g^*\omega)_j
		=(\nabla^*\nabla\omega)_j+\nabla_j\delta_g\omega
		-\operatorname{Ric}_j{}^a\omega_a.
	\]
	又 \(\operatorname{tr}_g\delta_g^*\omega=-\delta_g\omega\)，故加上 \(\mathrm d\operatorname{tr}_g\delta_g^*\omega\) 后中间项抵消，得到式\eqref{g3:eq:bianchi-killing-operator}。用 \(\partial_j\mapsto i\xi_j\) 取最高阶项，得到所列主符号；它与 \(\omega\) 的配对为 \(|\xi|^2|\omega|^2+|\langle\xi,\omega\rangle|^2>0\)。最后在闭流形上
	\[
		\langle\delta_g\delta_g^*\omega,\omega\rangle_{L^2(g)}
		=\|\delta_g^*\omega\|_{L^2(g)}^2,
	\]
	故核的描述成立。
\end{proof}

满足 \(\mathcal L_Xg=0\) 的向量场称为 \(g\) 的\term[Killing chang]{Killing 场}[Killing field]，它生成局部等距变换。记
\[
	\mathfrak k_0=\{\omega\in C^\infty(\mathrm T^*M):\delta_0^*\omega=0\}.
\]
由椭圆核的有限维性，\(\mathfrak k_0\) 有限维。Bianchi 条件 \(\beta_0(g-g_0)=0\) 也常用于 Einstein 方程；式\eqref{g3:eq:bianchi-killing-operator}说明它沿微分同胚方向产生 \(\nabla^*\nabla-\operatorname{Ric}\)。在正 Ricci 曲率下，不能仅凭这个表达式断言其核只有 Killing 场。散度条件给出非负算子 \(\delta_0\delta_0^*\)，核的几何意义则由上面的积分恒等式直接确定。

\begin{example}\label{g3:exm:bianchi-nonkilling-kernel}
	在两份截面曲率为 \(\kappa>0\) 的圆球乘积 \(S^2\times S^2\) 上，乘积度量满足 \(\operatorname{Ric}=\kappa g\)。设 \(f\) 为第一因子上的非零高度函数，拉回到乘积后有 \(\Delta f=2\kappa f\)、\(\operatorname{Hess}f=-\kappa f g_1\)。一形式的 Weitzenböck 公式给出
	\[
		\nabla^*\nabla\mathrm df
		=\mathrm d\Delta f-\operatorname{Ric}(\mathrm df,\cdot)
		=\kappa\,\mathrm df.
	\]
	因此 \(\beta\delta^*\mathrm df=0\)，但 \(\delta^*\mathrm df=\operatorname{Hess}f\ne0\)。这是 Bianchi 规范算子具有非 Killing 核的具体例子。
\end{example}

\subsection{规范化 Einstein 算子的主符号}
\label{g3:subsec:2502:02}
固定 \(\lambda\)，定义
\begin{equation}\label{g3:eq:einstein-conserved-gauge}
	\begin{aligned}
		\mathcal G_\lambda(g)&=\operatorname{Ric}_g-\tfrac12s_g g
			+\tfrac{n-2}{2}\lambda g,\\
		\omega(g)&=\delta_0(g-g_0),\\
		\mathcal F(g)&=\mathcal G_\lambda(g)+\delta_g^*\omega(g).
	\end{aligned}
\end{equation}
缩并 Bianchi 恒等式给出 \(\delta_g\mathcal G_\lambda(g)=0\)。另一方面，\(\mathcal G_\lambda(g)=0\) 的迹为 \((n-2)(n\lambda-s_g)/2=0\)，所以在 \(n\ge3\) 时它等价于 \(\operatorname{Ric}_g=\lambda g\)。以这个守恒张量代替 Ricci 张量，正是随后消去规范误差的原因。

\begin{theorem}\label{g3:thm:einstein-gauged-ellipticity}
	设 \(M^n\) 闭且光滑，\(n\ge3\)，\(g_0\) 为满足 \(\operatorname{Ric}_{g_0}=\lambda g_0\) 的光滑度量。取整数 \(s>n/2+3\)。式\eqref{g3:eq:einstein-conserved-gauge}定义光滑映射
	\[
		\mathcal F:\mathcal M^s(M)\longrightarrow
		H^{s-2}(S^2\mathrm T^*M).
	\]
	在 \(g_0\) 附近它是二阶准线性椭圆系统。其线性化
	\begin{equation}\label{g3:eq:einstein-gauged-linearization}
		Lh=\mathrm D\mathcal G_\lambda|_{g_0}h+\delta_0^*\delta_0h
	\end{equation}
	关于 \(L^2(g_0)\) 形式自伴。
\end{theorem}
\begin{proof}
	局部表达式中，逆度量是正定矩阵的光滑函数，曲率为 \(g^{-1}*\partial^2g+g^{-1}*g^{-1}*\partial g*\partial g\) 的有限和。规范项也至多包含二阶导数。\Cref{g3:prop:nonlinear-jet-operator}适用，因为 \(s-2>n/2\)，故映射光滑。由于 \(\omega(g_0)=0\)，微分规范项时只有 \(\delta_0^*\delta_0h\) 留下。

	令 \(\xi\) 为单位余向量，在 \(g_0\) 正交标架中取 \(\xi=e^1\)。上一节的 Ricci 线性化给出
	\[
		\sigma_{\operatorname{Ric}'}(\xi)h
		=\tfrac12\bigl(h+\xi\otimes\xi\operatorname{tr}h
		-\xi\otimes h(\xi^\sharp,\cdot)-h(\xi^\sharp,\cdot)\otimes\xi\bigr),
	\]
	而标量曲率的符号是 \(\operatorname{tr}h-h(\xi^\sharp,\xi^\sharp)\)。加上 \(\delta_0^*\delta_0\) 后，交叉项抵消。若 \(a,b\ge2\)、\(H=\sum_{c=2}^nh_{cc}\)，则
	\begin{equation}\label{g3:eq:einstein-gauge-block-symbol}
		(\sigma_Lh)_{11}=h_{11},\quad
		(\sigma_Lh)_{1a}=\tfrac12h_{1a},\quad
		(\sigma_Lh)_{ab}=\tfrac12(h_{ab}-H\delta_{ab}).
	\end{equation}
	最后一块的迹是 \(-(n-2)H/2\)。因此 \(\sigma_Lh=0\) 依次推出 \(h_{11}=h_{1a}=H=h_{ab}=0\)。任意非零 \(\xi\) 由齐次性得到可逆性。单位余球丛紧，最小奇异值有正下界；系数连续依赖 \(g\)，故附近度量仍椭圆。这个符号有负的迹方向，椭圆性并不等于能量正定。

	上一节的作用量 \(\mathcal A_\lambda(g)=\int_Ms_g\,\mu_g-(n-2)\lambda\operatorname{Vol}_g(M)\) 满足
	\(\mathrm D\mathcal A_\lambda(g)h=-\langle\mathcal G_\lambda(g),h\rangle_{L^2(g)}\)。在 \(g_0\) 处 \(\mathcal G_\lambda=0\)，故再微分时内积和体积的变化项都消失，Hessian 的对称性给出 \(\mathcal G_\lambda'\) 的形式自伴性。算子 \(\delta_0^*\delta_0\) 同样自伴，得到结论。
\end{proof}

\begin{theorem}\label{g3:thm:einstein-gauge-equivalence}
	设 \((M^n,g_0)\) 为闭光滑 Einstein 流形，\(n\ge3\)，\(\operatorname{Ric}_{g_0}=\lambda g_0\)。记 \(\delta_0=\delta_{g_0}\) 为负散度、\(\delta_g^*\) 为其对称梯度伴随，并令 \(\mathcal F(g)=\operatorname{Ric}_g-\tfrac12s_gg+\tfrac{n-2}2\lambda g+\delta_g^*\delta_0(g-g_0)\)。存在 \(g_0\) 的一个 \(C^1\) 邻域，使其中每个 \(C^3\) 度量 \(g\) 满足
	\[
		\mathcal F(g)=0
		\quad\Longleftrightarrow\quad
		\operatorname{Ric}_g=\lambda g\ \text{且}\ \delta_0(g-g_0)=0.
	\]
\end{theorem}
\begin{proof}
	只需证明正向。设 \(\omega=\delta_0(g-g_0)\)。对方程取 \(\delta_g\)，由守恒恒等式得到 \(\delta_g\delta_g^*\omega=0\)。与 \(\omega\) 配对并积分，得 \(\delta_g^*\omega=0\)，于是 \(\mathcal G_\lambda(g)=0\)，亦即 \(\operatorname{Ric}_g=\lambda g\)。

	还须证明代表条件本身。对任意 \(\eta\in\mathfrak k_0\)，有
	\(\langle\omega,\eta\rangle_0=\langle g-g_0,\delta_0^*\eta\rangle_0=0\)。在 \(\mathfrak k_0^\perp\) 上有估计
	\begin{equation}\label{g3:eq:killing-complement-first-order}
		\|\alpha\|_{H^1}\le C\|\delta_0^*\alpha\|_2.
	\end{equation}
	为证明它，先对具有单射符号的 \(\delta_0^*\) 冻结系数、用 Fourier 估计及有限覆盖，得 \(\|\alpha\|_{H^1}\le C(\|\delta_0^*\alpha\|_2+\|\alpha\|_2)\)。若不能去掉末项，取 \(\mathfrak k_0^\perp\) 中 \(H^1\) 范数为一而 \(\delta_0^*\alpha_j\to0\) 的序列。Rellich 紧性给出强 \(L^2\) 极限 \(\alpha\)；弱方程使它属于 \(\mathfrak k_0\)，正交性又使它为零。前述估计遂给 \(\alpha_j\to0\) 于 \(H^1\)，矛盾。

	两联络之差是由 \(g^{-1}\nabla^0g\) 给出的零阶张量，因此
	\[
		\|\delta_0^*\omega\|_2
		=\|(\delta_0^*-\delta_g^*)\omega\|_2
		\le C'\|g-g_0\|_{C^1}\|\omega\|_2.
	\]
	取邻域使 \(CC'\|g-g_0\|_{C^1}<1\)，式\eqref{g3:eq:killing-complement-first-order}迫使 \(\omega=0\)。反向代入定义即得。
\end{proof}

\subsection{切片所需的椭圆方程}
\label{g3:subsec:2502:03}
先把线性代表条件解出来。给定对称二张量 \(h\)，希望沿 \(h\mapsto h+\mathcal L_Xg_0\) 消去散度，因此要解
\begin{equation}\label{g3:eq:metric-linear-slice-equation}
	2\delta_0\delta_0^*X^\flat=-\delta_0h.
\end{equation}
右端自动与 \(\mathfrak k_0\) 正交。\Cref{g3:thm:elliptic-fredholm,g3:cor:elliptic-range-hk}因而给出唯一的核正交解，且 \(\|X\|_{H^{r+1}}\le C_r\|h\|_{H^r}\)。这也给出正交分解
\[
	H^r(S^2\mathrm T^*M)
	=\ker\delta_0\ \oplus\ \delta_0^*H^{r+1}(\mathrm T^*M),\qquad r\ge1.
\]
第二项不要求把 Killing 场重复计入参数。

非线性求解涉及变化源上的复合。为看清正则性，令
\(\phi_X(x)=\exp^{g_0}_xX(x)\)，其中 \(X\) 足够 \(C^1\) 小。它是邻近恒等的微分同胚：局部逆函数定理给局部可逆，紧性给覆盖，而与恒等的同伦使覆盖次数为一。

\begin{lemma}\label{g3:lem:metric-pullback-one-loss}
	设 \(M^n\) 闭且光滑，\(g_0\) 为光滑黎曼度量，整数 \(s>n/2+3\)。记 \(\phi_X(x)=\exp_x^{g_0}X(x)\)，\(\delta_0^*\) 为背景度量的对称梯度，\(\mathcal M^{s+1}(M)\) 为正定 \(H^{s+1}\) 对称二张量的开集。在 \(X=0\)、\(g=g_0\) 附近，映射
	\[
		(X,g)\longmapsto\phi_X^*g
	\]
	从 \(H^{s+1}(\mathrm TM)\times\mathcal M^{s+1}(M)\) 到 \(H^s(S^2\mathrm T^*M)\) 连续可微；其在 \((0,g_0)\) 处的微分是 \((X,h)\mapsto2\delta_0^*X^\flat+h\)。
\end{lemma}
\begin{proof}
	有限坐标覆盖上只需检查
	\((\phi^*g)_{ij}(x)=g_{ab}(\phi(x))\partial_i\phi^a\partial_j\phi^b\)。先说明复合估计：在 \(H^{s+1}\) 有界且 Jacobian 及其逆有界的微分同胚邻域中，\(f\mapsto f\circ\phi\) 在每个整数 \(0\le r\le s+1\) 的 \(H^r\) 上有统一界。\(r=0\) 由换元得出；从 \(r-1\) 到 \(r\) 对复合求一次导数，用 \(\mathrm D(f\circ\phi)=(\mathrm Df\circ\phi)\mathrm D\phi\) 及\cref{g3:lem:sobolev-multiple-products}的乘积估计归纳。对光滑 \(f\) 有关于 \(\phi\) 的连续性，一般 \(f\) 由光滑稠密性和统一界逼近，因此联合连续。

	局部坐标内，\(\phi\) 的增量 \(V\) 给出
	\[
		(f\circ(\phi+V)-f\circ\phi)
		=\int_0^1(\mathrm Df\circ(\phi+tV))V\,\mathrm dt.
	\]
	减去 \((\mathrm Df\circ\phi)V\)，联合连续性说明余项是 \(o(\|V\|_{H^{s+1}})\) 于 \(H^s\)。同一积分式与统一复合界还给出局部估计
	\[
		\|f\circ\phi-f\circ\psi\|_{H^s}
		\le C\|f\|_{H^{s+1}}\|\phi-\psi\|_{H^{s+1}}.
	\]
	它保证关于 \(f\) 的微分在算子范数中连续，并控制 \(f\) 与 \(\phi\) 同时变化时的交叉项；关于 \(\phi\) 的微分则由 \(\mathrm Df\circ\phi\) 在 \(H^s\) 中的连续性控制。故复合从两份 \(H^{s+1}\) 到 \(H^s\) 为 \(C^1\)。拉回表达式中的两个 \(\mathrm D\phi\) 属于 \(H^s\)，而 \(H^s\) 是乘法代数；代入纤维指数图，链式法则给出结论。微分的坐标表达式为 Lie 导数，故得到所列公式。
\end{proof}

以下正则性事实也将用于下一节的有限维约化。需要的是一般等秩椭圆系统，而不是只含标量主部的方程。

\begin{lemma}\label{g3:lem:elliptic-system-holder-bootstrap}
	设 \(u\) 是开集上的 \(C^{3,\alpha}\) 向量值函数，\(0<\alpha<1\)，满足
	\[
		A^{ij}(x,u,\mathrm Du)\partial_i\partial_ju
		=B(x,u,\mathrm Du)+f(x),
	\]
	其中 \(A,B,f\) 光滑，主符号 \(\sum A^{ij}(x,u,\mathrm Du)\xi_i\xi_j\) 在每个紧子集和单位 \(\xi\) 上可逆且逆一致有界。则 \(u\) 光滑。若主系数只依赖 \(x,u\)，结论对相同起始正则性同样成立；在固定紧子集上，每阶估计由相应低阶范数、椭圆性和数据界控制。
\end{lemma}
\begin{proof}
	先把第12.5节的内部 Schauder 证明扩展到矩阵主符号。设冻结算子为 \(P_0\)，其二阶符号为 \(p_2(\xi)\)。对二阶导数求解所用的乘子是 \(-\xi_k\xi_l p_2(\xi)^{-1}\)，它是零次齐次、在单位球面光滑的矩阵函数。分离常矩阵乘法项后，\cref{g3:thm:smooth-multiplier-lp}证明中的二进分解给出齐次核 \(K\) 的界 \(|K(x)|\le C|x|^{-n}\)、\(|\mathrm DK(x)|\le C|x|^{-n-1}\)，以及球面零均值所给的环域消去。对 \(f\in C^{0,\alpha}\)，在近区减去 \(f(x)\)，积分界为 \(C[f]_\alpha\int_0^r t^{\alpha-1}\,\mathrm dt\)；比较相距 \(r\) 的两点时，远区用核差，界为 \(Cr[f]_\alpha\int_r^1t^{\alpha-2}\,\mathrm dt\)。两项均为 \(O(r^\alpha)\)。取截断函数 \(\chi\)，把上述全空间估计用于 \(\chi v\)，利用 \(P_0(\chi v)=\chi P_0v+[P_0,\chi]v\)；交换子只含一阶导数，可用加权插值吸收。这就给出常系数系统的内部二阶 Hölder 估计。

	对变系数系统，冻结误差仍是矩阵乘积 \((A(x)-A(x_0))\mathrm D^2v\)。缩小球吸收最高阶半范数，并用\cref{g3:lem:weighted-holder-interpolation}吸收其余导数项，得到内部估计
	\[
		\|v\|_{C^{2,\alpha}(U')}
		\le C\bigl(\|Av\|_{C^{0,\alpha}(U)}+\|v\|_{C^0(U)}\bigr),
		\qquad U'\Subset U.
	\]
	此处只使用矩阵范数和逆符号界。对初始 \(C^2\) 解，\cref{g3:eq:holder-mollifier-commutator}的卷积交换子估计逐矩阵项适用，故估计包含正则性结论。对差商使用这个估计，正如\cref{g3:cor:schauder-higher-bootstrap}的证明，系数和右端为 \(C^{r,\alpha}\) 时得到 \(C^{r+2,\alpha}\) 正则性。

	回到本引理，初始 \(u\in C^{3,\alpha}\) 使 \(A(x,u,\mathrm Du)\) 和右端至少属于 \(C^{2,\alpha}\)。将这些已经确定的函数作为线性方程的系数，得到 \(u\in C^{4,\alpha}\)。系数和右端随之成为 \(C^{3,\alpha}\)，再得 \(C^{5,\alpha}\)。归纳完成证明。每一步的估计也给出所述依赖关系。
\end{proof}

\begin{theorem}\label{g3:thm:metric-divergence-slice}
	设 \(M^n\) 闭且光滑，\(g_0\) 为光滑黎曼度量，整数 \(s>n/2+3\)。记 \(\delta_0=\delta_{g_0}\)、\(\mathfrak k_0=\ker\delta_0^*\)，正交取 \(L^2(g_0)\) 内积；对充分小的 \(X\)，令 \(\phi_X(x)=\exp_x^{g_0}X(x)\)。每个充分接近 \(g_0\) 的 \(H^{s+1}\) 度量 \(g\)，均有唯一充分小的 \(X\in H^{s+1}(\mathrm TM)\)、\(X^\flat\perp\mathfrak k_0\)，使
	\[
		\delta_0(\phi_X^*g-g_0)=0.
	\]
	此时 \(\phi_X^*g\in H^s\)，\(X\) 对 \(g\) 连续可微。若 \(g\) 光滑，则 \(X\)、\(\phi_X\) 及所得代表都光滑。
\end{theorem}
\begin{proof}
	令 \(\Phi(X,g)=\delta_0(\phi_X^*g-g_0)\)。前一复合引理说明它取值于 \(H^{s-1}\) 且为 \(C^1\)。其像始终正交于 \(\mathfrak k_0\)，因为背景散度的像具有这个性质。在 \((0,g_0)\) 处，\(X\) 微分是 \(2\delta_0\delta_0^*X^\flat\)。\Cref{g3:prop:bianchi-gauge-symbol}与椭圆 Fredholm 选择表明它是
	\[
		H^{s+1}(\mathrm T^*M)\cap\mathfrak k_0^\perp
		\longrightarrow H^{s-1}(\mathrm T^*M)\cap\mathfrak k_0^\perp
	\]
	之间的有界同构。应用\cref{g3:cor:banach-implicit-parameters}，得到存在、唯一性及 \(C^1\) 依赖。

	若 \(g\) 光滑，写出 \(\Phi(X,g)=0\) 的坐标表达式，它是关于 \(X\) 的二阶准线性系统；最高阶系数光滑依赖 \(x,X,\mathrm DX\)，在 \(X=0,g=g_0\) 处为上述强椭圆符号。邻域缩小后保持椭圆。由 \(X\in H^{s+1}\) 得 \(X\in C^{3,\alpha}\) 对某个 \(\alpha>0\) 成立，故\cref{g3:lem:elliptic-system-holder-bootstrap}给出光滑性。
\end{proof}

这个证明把拉回的一阶微分损失明确放进了空间选择：输入度量为 \(H^{s+1}\)，所得代表先在 \(H^s\) 中控制。它没有把微分同胚作用当作同一 Sobolev 阶上的光滑作用；光滑度量的几何结论则由椭圆正则性恢复。

\subsection{无穷小等距与剩余对称}
\label{g3:subsec:2502:04}
令 \(I_0=\operatorname{Isom}(M,g_0)\)。它保持散度条件，也保持规范化方程：对 \(a\in I_0\)，
\[
	\delta_0(a^*h)=a^*(\delta_0h),\qquad
	\mathcal F(a^*g)=a^*\mathcal F(g).
\]
因而规范固定后仍须除去 \(I_0\) 的作用。其 Lie 代数为 Killing 场，但离散等距变换同样可能起作用。

\begin{proposition}\label{g3:prop:metric-slice-residual-isometries}
	设 \(M\) 闭且连通，\(g_0\) 光滑。记 \(\delta_0=\delta_{g_0}\) 为负散度，\(I_0=\operatorname{Isom}(M,g_0)\)。光滑度量在 \(g_0\) 附近满足 \(\delta_0(g-g_0)=0\) 时，模所有微分同胚的局部等价关系恰由 \(I_0\) 的拉回作用给出。这里“附近”取光滑拓扑，并可相应缩小上一切片定理的邻域。
\end{proposition}
\begin{proof}
	先考察恒等附近的微分同胚。Killing 场的流保持 \(g_0\)，由闭性对所有小时间存在。选 \(\mathfrak k_0\) 的一组基，按次序复合这些流，得到有限维参数 \(a(t)\)。映射 \((X,t)\mapsto a(t)\circ\phi_X\)，其中 \(X^\flat\perp\mathfrak k_0\)，在 \((0,0)\) 的微分是核正交向量场与 Killing 场之和，故为同构。这里外层流光滑且仅有有限个参数，映射在 \(H^{s+1}\) 微分同胚图中光滑。逆函数定理说明每个充分靠近恒等的微分同胚都有这种表示。

	若 \(g_1,g_2\) 均在散度切片内，且 \(g_2=(a(t)\circ\phi_X)^*g_1=\phi_X^*(a(t)^*g_1)\)，则 \(a(t)^*g_1\) 也在切片内。切片定理对这一输入的唯一解是 \(X=0\)，因此该等价来自等距变换。

	最后排除远离 \(I_0\) 的额外等价。若有 \(g_j,h_j\to g_0\) 光滑，\(\psi_j^*g_j=h_j\)，则拉回恒等式给出 \(\mathrm D\psi_j\) 及其逆的一致界。局部等距映射保持联络，因此在坐标中
	\[
		\partial_i\partial_j\psi_j^a
		=\Gamma(h_j)^k_{ij}\partial_k\psi_j^a
		-\Gamma(g_j)^a_{bc}(\psi_j)\partial_i\psi_j^b\partial_j\psi_j^c.
	\]
	有限坐标覆盖与此式依次控制全部高阶导数；逆映射有相同估计。由 Arzelà--Ascoli 及对角取列，\(\psi_j\) 和逆映射光滑收敛到互逆的映射 \(\psi,\psi^{-1}\)，并且 \(\psi^*g_0=g_0\)。故 \(\psi\in I_0\)，\(\psi^{-1}\circ\psi_j\to\operatorname{id}\)。应用上一段，便得到所需结论；否则一列违反结论的等价会产生矛盾。
\end{proof}

同一紧性论证表明 \(I_0\) 紧。恒等附近的分解还说明其中靠近恒等的元素全部由 Killing 场的有限维流参数给出：对 \(g_1=g_2=g_0\) 再用切片唯一性即可。于是 \(I_0\) 是有限维紧 Lie 群；这也说明为何它的离散部分不能由核空间计算替代。上面的论证是本章使用的散度切片版本；度量空间切片的一般框架见\cite[第7节，定理7.4，第33--34页]{Ebin1970Metrics}。
% 纲要条目排为编号小节标题，并进入目录与书签；说明段仍保留于纲要注释中。
